\chapter*{Introduction}
\addcontentsline{toc}{chapter}{Introduction}
\markboth{Introduction}{Introduction}
\label{ch:prologo}

\epigraph{``Omnia, Lucili, aliena sunt, tempus tantum nostrum
est.''}{Lucius Annaeus Seneca, \emph{Epistulae morales ad
Lucilium}, I, 1}

$\pi$Tantum is a web application that assists in the
composition of Italian-school timetables. It originates from
ideas discussed and developed by Fernando Gargiulo, Giovanni
Borghi, Matteo Mariani and Stefano Bertozzi.

The school timetabling problem consists of assigning, to each
\emph{cattedra} -- the pair formed by a teacher and a class for
a specific subject -- a set of weekly hours distributed over
days and time slots, such that all the following constraints
are honoured: teacher availability, curricular hour coverage,
classroom and equipment compatibility, and the constraints that
derive from the Italian national collective contract for the
school sector. This is a combinatorial problem classified as
NP-hard, which means that as a school grows in size there is no
fast algorithm capable of finding an exact optimal solution.

$\pi$Tantum delegates the actual solving to the CP-SAT solver in
Google OR-Tools, chosen for its maturity in constraint
programming applied to scheduling problems. To scale beyond the
size that the solver would handle directly, the system
decomposes the problem into smaller sub-problems via four
alternative strategies: spectral decomposition of the bipartite
class-teacher graph, temporal decomposition by day of the week,
curriculum-based decomposition, and balanced $k$-way METIS
partitioning. On top of these, a cascade of metaheuristics
(Large Neighborhood Search, Adaptive LNS, Simulated Annealing,
Tabu Search, Iterated Local Search and Variable Neighborhood
Search) improves the quality of the solution found.

The application backend is written in Python on FastAPI; data
persistence is delegated to SQLite. The user interface is
based on SvelteKit (in Italian for the production deployment),
organised into sixteen tabs that cover the whole lifecycle of a
timetable: registry import, constraint configuration, generation,
manual editing with real-time feasibility checks, and
day-to-day handling of absences and substitutions.

The following chapters describe the problem formally, the
entities involved, the user interface, and each of the
algorithms that $\pi$Tantum employs to compute the timetable.
The exposition is intentionally multi-layered: a non-technical
reader can follow the narrative thread without engaging the
mathematical formulas, while a developer will find the
implementation details, code references, and source-file
pointers necessary to extend the system.

\bigskip

\noindent
This English edition is a parallel translation of the Italian
manual (\texttt{docs/manual.tex}). The Italian edition
remains the primary reference; this English edition is
intended for international readers and contributors. Where
Italian-school context is necessary (e.g.\ the legal frame of
the \emph{contratto collettivo nazionale}, the role of
\emph{insegnante di sostegno}, or the meaning of
\emph{indirizzo di studio}), explanatory notes are added.
